\section{Introduction}
\label{sec:introduction}

Large language models are increasingly used as research assistants in finance,
from domain-specific models~\citep{wu2023bloomberggpt} to studies that ask
whether model judgments predict returns~\citep{lopezlira2023chatgpt}. A common
design gives the model tools that read market data and compute statistics,
following the reasoning-and-acting pattern of tool-using agents~\citep{yao2023react,schick2023toolformer}.
A copilot of this kind is only useful for research if three properties hold,
and none of them can be checked by reading the final answer.

\paragraph{Point-in-time discipline.}
A question posed ``as of'' a date $t$ must be answered with information
available at $t$. Look-ahead bias is the classic failure of backtests, and
language models add a second channel: facts memorized during
training~\citep{glasserman2023lookahead}. A tool layer that silently clamps a
request for a later date to $t$ hides the attempt. An answer built from future
rows can also look correct to an evaluator who scores it against hindsight.

\paragraph{Numeric grounding.}
Every number in an answer should be traceable to a recorded computation.
Fluent but unsupported output is a well-documented failure of language
generation~\citep{ji2023hallucination}. In quantitative work a plausible
volatility that no tool computed is worse than no answer.

\paragraph{Action safety.}
A copilot that can reach an order-entry path can be talked into trading,
including by instructions that claim prior approval. Research governance
therefore requires that agents may propose bounded actions but never execute
them without explicit human confirmation.

\medskip
This paper asks whether a tool-using LLM copilot can answer quantitative
market questions with verifiable grounding while respecting point-in-time and
execution constraints, and how to measure that. Our position is that each
property has to be made \emph{observable by construction} before it can be
measured. The contributions are:

\begin{enumerate}
  \item \textbf{A governed copilot architecture} (Section~\ref{sec:system}).
  Data are read only through a confirmed-only point-in-time view. The as-of
  clock refuses rather than clamps later dates, so attempts can be counted. A
  policy gate checks every tool call before any handler runs and cannot be
  configured to allow execution. Order ideas become inert proposals that await
  human approval.
  \item \textbf{A point-in-time benchmark with trap categories}
  (Section~\ref{sec:benchmark}). The \SuiteTasks{} deterministic tasks
  include cutoff traps worded exactly like answerable questions, trade requests
  with false authority claims, and unknown or not-yet-listed symbols. Ground
  truth comes from an independent reference implementation, and hindsight
  values reveal answers built from future data.
  \item \textbf{A rounding-aware numeric grounding metric}
  (Section~\ref{sec:grounding}) that attributes each numeric claim to a
  citation and accepts it only if it is a correct rounding, with the right
  sign, of an output of the cited result that the text actually refers to.
  \item \textbf{An audit and provenance design}
  (Section~\ref{sec:provenance}). Content-addressed result identifiers are
  derived from the exact rows used, and a hash-chained log whose head is
  recorded in the committed run summary makes the record of every model call
  and tool attempt tamper-evident.
\end{enumerate}

\paragraph{Status of the evidence.}
This draft reports only \emph{harness validation}: scripted baselines and
synthetic claims on synthetic data, which show that the instruments register
what they claim to register (Section~\ref{sec:results}). The same harness
drives a Claude backend, and the protocol for language-model runs, with its
hypotheses and power analysis, is fixed in Section~\ref{sec:experiments}.
Those runs have not been executed.
\pending{language-model results for RQ1--RQ4}
